% Generated draft: requires mathematical and source review.
\documentclass[11pt]{article}
\usepackage[margin=25mm]{geometry}
\usepackage{fontspec}
\setmainfont{Latin Modern Roman}
\usepackage{amsmath,amssymb,mathtools,unicode-math}
\setmathfont{Latin Modern Math}
\usepackage{enumitem,needspace,xurl,hyperref}
\setlength{\parindent}{0pt}
\setlength{\parskip}{5pt}
\setlength{\emergencystretch}{3em}
\newcommand{\N}{\mathbb{N}}
\newcommand{\Z}{\mathbb{Z}}
\newcommand{\Q}{\mathbb{Q}}
\newcommand{\R}{\mathbb{R}}
\newcommand{\C}{\mathbb{C}}
\newcommand{\F}{\mathbb{F}}
\newcommand{\A}{\mathbb{A}}
\newcommand{\PP}{\mathbb{P}}
\newcommand{\E}{\mathbb{E}}
\newcommand{\eps}{\varepsilon}
\newcommand{\dd}{\,\mathrm{d}}
\begin{document}
\section*{UnsolvedMath ID: 9400094}
\textbf{Source problem code:} catalogue entry
\subsection{Definitions and mathematical statement}
Let \(E/\mathbb F_q\) be an elliptic curve, let \(P\in E(\mathbb F_q)\) have order \(r\), and let \(Q\in\langle P\rangle\). The elliptic-curve discrete-logarithm problem asks for \(a\in\mathbb Z/r\mathbb Z\) such that
\[
 [a]P=Q.
\]
The source asks whether it can be solved in subexponential time. On the usual classical, general-curve reading, seek an algorithm applicable uniformly to the intended unrestricted input family, with running time subexponential in the bit-length parameter, rather than only for a special curve family.

For example, a broad meaning with parameter \(L\) is \(\exp(o(L))\), where \(L\) includes the field, curve and point encoding lengths. A commonly used narrower subexponential scale is
\[
 L_q(\alpha,c)=\exp\bigl((c+o(1))(\log q)^\alpha(\log\log q)^{1-\alpha}\bigr),
 \qquad0<\alpha<1.
\]
The excerpt does not choose a scale, expected versus worst-case bound, success convention, field regime, or classical versus quantum model. Those must be fixed before a single complexity assertion is stated.

The summary reports no general classical subexponential method and mentions Pollard rho's square-root generic complexity. Its informal \(O(\sqrt p)\) is naturally expressed as \(O(\sqrt r)\) group operations for subgroup order \(r\), not automatically the field characteristic. Generic-group bounds do not establish lower bounds for all algorithms using curve representations. The summary also reports faster special-family methods, naming anomalous, supersingular and specially structured extension-field curves; these do not constitute a uniform general-curve algorithm.

The pinned Wikipedia open label and the later hardness assessment are source reports, not a proof of computational impossibility or an independently verified current status.
\subsection{Short English statement}
Can the general elliptic-curve discrete logarithm be solved by a classical algorithm in time subexponential in input size? The source's intended computation model and time scale need confirmation; square-root generic methods and special-curve algorithms are distinct from the general target.
\subsection{Sources}
\begin{itemize}
\item Local registry: \url{data/unsolvedmath/problems.json}, permanent UnsolvedMath ID \(9400094\), source code catalogue entry; supplied statement, background, and research summary, updated 2026-08-06.
\item Wikipedia contributors, \emph{List of unsolved problems in mathematics}, Number theory, computational number theory bullet 94, wikitext line 116, pinned revision 1366636767: \url{https://en.wikipedia.org/w/index.php?title=List_of_unsolved_problems_in_mathematics&oldid=1366636767#Number_theory}. MediaWiki-source extraction accessed 2026-07-29; primary-source audit and final semantic deduplication remain pending in the provenance.
\item The summary names Pollard rho and the Wikipedia ``Elliptic-curve discrete logarithm'' page, with no exact algorithm or special-case publication citations.
\end{itemize}
\end{document}
